\documentclass[11pt, a4paper]{exam}
\usepackage[a4paper,top=2.3cm,bottom=2.8cm,left=2cm,right=2cm,headheight=15pt,headsep=0.6cm,footskip=1.2cm]{geometry}
\usepackage[T1]{fontenc}
\usepackage[utf8]{inputenc}
\usepackage{lmodern}
\usepackage{textcomp}
\usepackage[french]{babel}
\usepackage{xcolor}
\usepackage{listings}
\usepackage{graphicx}
\usepackage{array}
\usepackage{enumitem}
\usepackage{amsmath}
\usepackage{tikz}
\usetikzlibrary{positioning,shapes.multipart,arrows.meta}
\usepackage[most]{tcolorbox}

\definecolor{codegreen}{rgb}{0,0.6,0}
\definecolor{codegray}{rgb}{0.5,0.5,0.5}
\definecolor{codepurple}{rgb}{0.58,0,0.82}
\definecolor{backcolour}{rgb}{0.95,0.95,0.92}

\lstdefinestyle{mystyle}{
    backgroundcolor=\color{backcolour},
    commentstyle=\color{codegreen},
    keywordstyle=\color{magenta},
    numberstyle=\tiny\color{codegray},
    stringstyle=\color{codepurple},
    basicstyle=\ttfamily\footnotesize,
    breakatwhitespace=false,
    breaklines=true,
    postbreak=\mbox{\textcolor{codegray}{$\hookrightarrow$}\space},
    captionpos=b,
    keepspaces=true,
    columns=fullflexible,
    numbers=left,
    numbersep=5pt,
    showspaces=false,
    showstringspaces=false,
    showtabs=false,
    upquote=true,
    tabsize=4,
    xleftmargin=1.5em,
    framexleftmargin=0.3em,
    aboveskip=0.6em,
    belowskip=0.6em,
    extendedchars=true,
    literate={é}{{\'e}}1 {è}{{\`e}}1 {ê}{{\^e}}1 {ë}{{\"e}}1 {à}{{\`a}}1 {â}{{\^a}}1
             {ç}{{\c c}}1 {ù}{{\`u}}1 {û}{{\^u}}1 {ô}{{\^o}}1 {î}{{\^i}}1 {ï}{{\"i}}1
             {É}{{\'E}}1 {È}{{\`E}}1 {À}{{\`A}}1 {Ç}{{\c C}}1 {œ}{{\oe}}1
             {«}{{\guillemotleft}}1 {»}{{\guillemotright}}1 {→}{{$\rightarrow$}}1
}
\lstset{style=mystyle, language=Python}
\lstdefinestyle{sqlstyle}{style=mystyle, language=SQL, morekeywords={LIMIT,REFERENCES,IN,TEXT,FLOAT,DATE,VARCHAR,CHAR,REPLACE}, numbers=none, xleftmargin=0.3em, frame=lines}

% Code en ligne
\newcommand{\code}[1]{\texttt{#1}}

% Titre d'une question : numéro et intitulé, barème à droite
\newcommand{\Qt}[2]{\subsection*{#1\hfill{\normalfont\small #2}}}
% Rappel de la consigne en une ligne
\newcommand{\rappel}[1]{\noindent\textit{Ce qui est demandé :} #1\par\smallskip}

% La réponse attendue, mise en évidence
\newtcolorbox{reponse}[1][Réponse]{enhanced, breakable, colback=black!5, colframe=black, boxrule=0.9pt, arc=1mm,
  title={#1}, fonttitle=\bfseries\small, coltitle=black, colbacktitle=black!15, left=2mm, right=2mm, top=1mm, bottom=1mm,
  before skip=0.5em, after skip=0.7em}
% Erreur fréquente
\newtcolorbox{erreur}{enhanced, breakable, colback=white, colframe=black, boxrule=0.5pt, leftrule=3pt, arc=0mm,
  title={Erreur fréquente}, fonttitle=\bfseries\small, coltitle=black, colbacktitle=white, left=2mm, right=2mm, top=1mm, bottom=1mm,
  before skip=0.5em, after skip=0.7em}
% Encadré « À retenir »
\newtcolorbox{retenir}{enhanced, breakable, colback=white, colframe=black, boxrule=0.5pt, arc=1mm, boxsep=1mm,
  borderline={0.4pt}{1.2pt}{black, dashed},
  title={À retenir}, fonttitle=\bfseries\small, coltitle=black, colbacktitle=white, left=2mm, right=2mm, top=1mm, bottom=1mm,
  before skip=0.5em, after skip=0.7em}

\newcommand{\methode}{\noindent\textbf{Méthode.}\ }
\newcommand{\explication}{\noindent\textbf{Explication.}\ }

\setlist{itemsep=0.15em, topsep=0.3em}
\renewcommand{\arraystretch}{1.25}
\setlength{\parindent}{0pt}
\setlength{\parskip}{0.35em}

% Piles dessinées verticalement (sommet en haut) ; \suppr barre un élément supprimé
\newcommand{\suppr}[1]{\ooalign{#1\cr\hidewidth\rule[0.55ex]{1em}{1pt}\hidewidth\cr}}
\newcommand{\pilev}[2]{\begin{tikzpicture}[baseline=(titre.base)]
  \node[font=\scriptsize, align=center, text width=22mm] (titre) at (0,0) {#1};
  \foreach \v [count=\k] in {#2} {\node[draw, minimum width=8mm, minimum height=5mm, inner sep=0pt, font=\small] at (0,-0.25-\k*0.5) {\v};}
\end{tikzpicture}}
\setlength{\emergencystretch}{3em}
% Remarque sur l'énoncé (erreur ou ambiguïté du sujet)
\newtcolorbox{remarque}{enhanced, breakable, colback=white, colframe=black, boxrule=0.5pt, arc=0mm,
  borderline={0.4pt}{1.2pt}{black, dotted},
  title={Remarque sur l'énoncé}, fonttitle=\bfseries\small, coltitle=black, colbacktitle=black!8, left=2mm, right=2mm, top=1mm, bottom=1mm,
  before skip=0.5em, after skip=0.7em}
\newcommand{\sql}[1]{\textsc{#1}}
\lstset{style=sqlstyle}
\newcommand{\resultat}[1]{\par\smallskip\noindent\textit{Résultat :}\par\nopagebreak\begin{center}\small #1\end{center}}

\title{Corrigé — CC Term NSI - SQL}
\author{Babash - NSI Tani Malandi et Johanne Bakalara}
\date{Contrôle de novembre 2024 (1 h 30, sur 20 points)}

\begin{document}
\pagestyle{headandfoot}
\firstpageheader{}{}{}
\runningheader{\small Corrigé}{}{\small SQL, location de matériel (d'après le bac 2023)}
\runningheadrule
\firstpagefooter{}{\footnotesize Babash - NSI Tani Malandi et Johanne Bakalara \textperiodcentered{} CC BY-SA 4.0 \textperiodcentered{} page \thepage}{}
\runningfooter{}{\footnotesize Babash - NSI Tani Malandi et Johanne Bakalara \textperiodcentered{} CC BY-SA 4.0 \textperiodcentered{} page \thepage}{}

\maketitle
\thispagestyle{headandfoot}
\begin{center}\footnotesize\itshape D'après l'exercice 2 du sujet de NSI du baccalauréat 2023 (jour 2, Polynésie, code 23-NSIJ2PO1).\end{center}

\begin{tcolorbox}[colback=white, colframe=black, boxrule=0.5pt, arc=1mm, left=2mm, right=2mm]
\small\textbf{Comment lire ce corrigé.} Pour chaque question : ce qui est demandé, puis la réponse attendue (encadrée), puis l'explication pas à pas. Essaie d'abord seul·e, puis compare ta démarche, pas seulement ton résultat.
Toutes les requêtes ont été exécutées avec SQLite sur une base qui contient exactement les tables du sujet ; le résultat obtenu est donné sous chaque requête.
Barème du sujet : 1 point par question, 2 points par requête à rédiger (20 points).
\end{tcolorbox}

\begin{retenir}
\begin{itemize}
\item Une \textbf{relation} est une table ; un \textbf{attribut} est une colonne ; un \textbf{enregistrement} (ou n-uplet) est une ligne.
\item Une \textbf{clé étrangère} (notée \#) est un attribut qui reprend la clé primaire d'une autre table : \code{Reservation.id\_membre} renvoie à \code{Membre.id\_membre}. C'est elle qui permet les \textbf{jointures}.
\item Une requête de lecture se construit toujours dans l'ordre : \sql{SELECT} (quelles colonnes ?) \sql{FROM} (quelle table ?) \sql{JOIN} \ldots{} \sql{ON} (quelles tables relier, et par quelle égalité ?) \sql{WHERE} (quelles lignes garder ?).
\end{itemize}
\end{retenir}

\section*{Partie 1 -- Lire les tables (sans requête)}
\Qt{Question 1}{4 points}

\textbf{a)} \rappel{la définition d'une clé primaire.}
\begin{reponse}
Une \textbf{clé primaire} est un attribut (ou un groupe d'attributs) dont la valeur identifie \textbf{de façon unique} chaque enregistrement de la table : deux lignes ne peuvent pas avoir la même valeur de clé primaire, et cette valeur ne peut pas être vide.
\end{reponse}
\explication Dans \code{Membre}, deux personnes pourraient avoir le même nom (deux « Dupont ») : le nom ne peut pas servir de clé. Le numéro \code{id\_membre}, lui, est différent pour chaque membre.

\textbf{b)} \rappel{les membres qui proposent un appareil à raclette.}
\begin{reponse}
Mohamed Ali et Fernando Alonso.
\end{reponse}
\methode On passe d'une table à l'autre en suivant les identifiants :
\begin{enumerate}
\item Dans \code{Objet}, l'appareil à raclette a l'\code{id\_objet} 4.
\item Dans \code{Possede}, les lignes avec \code{id\_objet} = 4 sont (1, 4) et (2, 4) : ce sont les membres 1 et 2.
\item Dans \code{Membre}, le membre 1 est Mohamed Ali et le membre 2 est Fernando Alonso.
\end{enumerate}

\textbf{c)} \rappel{le membre qui ne propose aucun objet.}
\begin{reponse}
Harry Kane.
\end{reponse}
\explication Les \code{id\_membre} présents dans \code{Possede} sont 1, 2, 3 et 4. Le membre 5 n'y apparaît pas : c'est Harry Kane. (Il a loué un objet, réservation 1, mais il n'en propose pas.)

\textbf{d)} \rappel{le schéma relationnel de \code{Reservation}, avec le domaine de chaque attribut.}
\begin{reponse}
Reservation(\underline{id\_reservation} : INT, \#id\_objet : INT, \#id\_membre : INT, date\_location : DATE, date\_retour : DATE)
\end{reponse}
\explication On suit le modèle donné pour \code{Objet} et \code{Possede} : la clé primaire est \underline{soulignée}, les clés étrangères sont précédées de \#, et chaque attribut est suivi de son \textbf{domaine} (le type de ses valeurs). Les domaines se lisent sur la figure du schéma relationnel. On peut aussi préciser NOT NULL pour la clé primaire, comme dans le schéma d'\code{Objet}.

\section*{Partie 2 -- Lire des requêtes}
\Qt{Question 2}{3 points}

\textbf{a)} \rappel{le résultat de \code{SELECT nom, prenom FROM Membre WHERE cp = "69003";}}
\begin{reponse}
\begin{tabular}{|l|l|}\hline
nom & prenom \\\hline
Dupont & Antoine \\
Kane & Harry \\\hline
\end{tabular}
\end{reponse}
\explication La requête garde les lignes de \code{Membre} dont le code postal est \code{"69003"} (membres 3 et 5), et n'affiche que les colonnes \code{nom} et \code{prenom}, dans cet ordre.

\textbf{b)} \rappel{le résultat de \code{SELECT MAX(tarif) FROM Objet;}}
\begin{reponse}
Une seule valeur : \textbf{20}.
\end{reponse}
\explication \sql{MAX} est une \textbf{fonction d'agrégation} : elle regroupe toutes les lignes en une seule valeur, ici le plus grand tarif (celui du nettoyeur haute pression). La requête ne donne pas le nom de l'objet, seulement le nombre.

\textbf{c)} \rappel{expliquer ce que fait \code{SELECT * FROM Reservation WHERE date\_location LIKE "2022-\%";}}
\begin{reponse}
La requête affiche toutes les colonnes (\code{*}) des réservations dont la date de location commence par « 2022- », c'est-à-dire les réservations commencées en 2022. Ici, ce sont les 6 réservations de la table.
\end{reponse}
\explication Avec \sql{LIKE}, le symbole \code{\%} remplace n'importe quelle suite de caractères (même vide). \code{"2022-\%"} veut donc dire « commence par 2022- ». Les dates s'écrivent sous la forme année-mois-jour : le début donne l'année.

\section*{Partie 3 -- Écrire des requêtes}
\Qt{Question 3}{6 points}

\textbf{a)} \rappel{le tarif de location d'une scie circulaire. \hfill\textit{(2 points)}}
\begin{reponse}
\begin{lstlisting}
SELECT tarif FROM Objet WHERE description = "Scie circulaire";
\end{lstlisting}
\end{reponse}
\resultat{\begin{tabular}{|l|}\hline tarif \\\hline 15 \\\hline\end{tabular}}
\methode
\begin{enumerate}
\item Quelle colonne ? le tarif : \code{SELECT tarif}.
\item Dans quelle table ? \code{Objet}.
\item Quelles lignes ? celle dont la description est « Scie circulaire » : \code{WHERE description = "Scie circulaire"}.
\end{enumerate}

\begin{erreur}
Écrire le texte sans guillemets (\code{WHERE description = Scie circulaire}) : SQL le prend pour un nom de colonne et signale une erreur. Une valeur de texte s'écrit toujours entre guillemets.
\end{erreur}

\textbf{b)} \rappel{les objets dont la description contient « appareil », en minuscules. \hfill\textit{(2 points)}}
\begin{reponse}
\begin{lstlisting}
SELECT * FROM Objet WHERE description LIKE "%appareil%";
\end{lstlisting}
\end{reponse}
\resultat{\begin{tabular}{|c|l|c|}\hline id\_objet & description & tarif \\\hline 4 & Appareil à raclette & 10 \\ 6 & Appareil à gaufre & 10 \\\hline\end{tabular}}
\explication « Contient » se traduit par un \code{\%} de chaque côté : n'importe quoi avant, le mot, n'importe quoi après. Avec un seul \code{\%} à la fin (\code{"appareil\%"}), on chercherait les descriptions qui \textbf{commencent} par « appareil ».

\begin{remarque}
Dans la table, les descriptions s'écrivent « \textbf{A}ppareil », avec une majuscule. Le résultat dépend donc du logiciel :
\begin{itemize}
\item avec SQLite (le logiciel le plus utilisé en NSI), \sql{LIKE} ne tient pas compte des majuscules : la requête trouve les objets 4 et 6 (résultat ci-dessus) ;
\item avec un logiciel où \sql{LIKE} tient compte des majuscules (comme PostgreSQL), elle ne trouve \textbf{aucune} ligne.
\end{itemize}
La requête attendue est celle de l'encadré « Réponse ». Pour être sûr·e de trouver « Appareil » dans tous les cas, on peut écrire \code{WHERE LOWER(description) LIKE "\%appareil\%"} (\code{LOWER} met le texte en minuscules), mais cette fonction n'est pas au programme.
\end{remarque}

\textbf{c)} \rappel{passer le tarif du nettoyeur haute pression de 20 \texteuro{} à 15 \texteuro{} par jour. \hfill\textit{(2 points)}}
\begin{reponse}
\begin{lstlisting}
UPDATE Objet SET tarif = 15 WHERE description = "Nettoyeur haute pression";
\end{lstlisting}
\end{reponse}
\explication \sql{UPDATE} modifie des lignes existantes : \sql{UPDATE} \textit{table} \sql{SET} \textit{colonne} = \textit{nouvelle valeur} \sql{WHERE} \textit{condition}. La requête modifie une ligne ; après elle, l'objet 1 a le tarif 15. On peut aussi écrire \code{WHERE id\_objet = 1}.

\begin{remarque}
L'énoncé parle d'un « nettoyeur \textbf{à} haute pression », mais la table contient « Nettoyeur haute pression ». Avec \code{=}, il faut recopier exactement le texte de la table : avec le « à », la requête ne modifierait aucune ligne.
\end{remarque}

\begin{erreur}
Oublier le \sql{WHERE} : \code{UPDATE Objet SET tarif = 15;} met le tarif 15 à \textbf{tous} les objets.
\end{erreur}

\section*{Partie 4 -- Ajouter des données}
\Qt{Question 4}{3 points}

\textbf{a)} \rappel{ajouter Wendie Renard (Villeurbanne, 69100) avec l'\code{id\_membre} 6. \hfill\textit{(2 points)}}
\begin{reponse}
\begin{lstlisting}
INSERT INTO Membre VALUES (6, "Renard", "Wendie", "69100");
\end{lstlisting}
\end{reponse}
\explication \sql{INSERT INTO} \textit{table} \sql{VALUES} (\ldots) ajoute une ligne. Les valeurs se donnent \textbf{dans l'ordre des colonnes} de la table : \code{id\_membre}, \code{nom}, \code{prenom}, \code{cp}. La ville n'est pas une colonne de \code{Membre} : on ne l'écrit pas. Le code postal est un texte (TEXT), d'où les guillemets.

\begin{erreur}
Inverser le nom et le prénom (\code{"Wendie", "Renard"}) : la requête fonctionne, mais la donnée est fausse. Pour éviter cette erreur, on peut nommer les colonnes : \code{INSERT INTO Membre (id\_membre, nom, prenom, cp) VALUES (6, "Renard", "Wendie", "69100");}
\end{erreur}

\textbf{b)} \rappel{inscrire Camille Martin (Bouéni, 97620), qui loue le 2024-11-28 l'appareil à gaufre et la scie circulaire, rendus une semaine plus tard. \hfill\textit{(1 point)}}
\begin{reponse}
\begin{lstlisting}
INSERT INTO Membre VALUES (7, "Martin", "Camille", "97620");
INSERT INTO Reservation VALUES (7, 6, 7, "2024-11-28", "2024-12-05");
INSERT INTO Reservation VALUES (8, 5, 7, "2024-11-28", "2024-12-05");
\end{lstlisting}
\end{reponse}
\methode
\begin{enumerate}
\item Le nouveau membre a besoin d'un \code{id\_membre} libre : 1 à 6 sont pris (le 6 par la question a), on prend 7.
\item Chaque objet loué est une \textbf{réservation} : il y a deux objets, donc deux lignes dans \code{Reservation}, avec des \code{id\_reservation} libres (7 et 8, car 1 à 6 sont pris).
\item On retrouve les numéros des objets dans la table \code{Objet} : appareil à gaufre = 6, scie circulaire = 5.
\item Une semaine après le 28 novembre 2024, c'est le 5 décembre 2024 (novembre a 30 jours : 28 + 7 = 35, et 35 $-$ 30 = 5).
\item L'ordre des colonnes de \code{Reservation} est : \code{id\_reservation}, \code{id\_objet}, \code{id\_membre}, \code{date\_location}, \code{date\_retour}.
\end{enumerate}
Après ces trois requêtes, \code{SELECT * FROM Reservation WHERE id\_membre = 7;} donne :
\begin{center}\small\begin{tabular}{|c|c|c|c|c|}\hline id\_reservation & id\_objet & id\_membre & date\_location & date\_retour \\\hline 7 & 6 & 7 & 2024-11-28 & 2024-12-05 \\ 8 & 5 & 7 & 2024-11-28 & 2024-12-05 \\\hline\end{tabular}\end{center}

\begin{erreur}
Ajouter une ligne dans \code{Possede} : cette table dit qui \textbf{propose} un objet à la location. Camille Martin \textbf{loue} des objets, ce sont des réservations. Autre erreur : réutiliser un \code{id\_reservation} déjà pris (1 à 6) ; la clé primaire doit être unique, la base refuse la ligne.
\end{erreur}

\begin{remarque}
Si tu n'as pas fait la question a, Camille Martin peut recevoir l'\code{id\_membre} 6 : la réponse est alors juste avec 6 à la place de 7 dans les trois requêtes.
\end{remarque}

\section*{Partie 5 -- Jointures}
\Qt{Question 5}{4 points}

\begin{retenir}
Une \textbf{jointure} relie deux tables en associant les lignes qui ont la même valeur pour la clé étrangère et la clé primaire :
\code{FROM Reservation JOIN Membre ON Reservation.id\_membre = Membre.id\_membre}.
On écrit \textit{table}.\textit{colonne} quand une colonne porte le même nom dans les deux tables.
\end{retenir}

\textbf{a)} \rappel{compter les réservations de Fernando Alonso, sans connaître son numéro. \hfill\textit{(2 points)}}
\begin{reponse}
\begin{lstlisting}
SELECT COUNT(*)
FROM Reservation
JOIN Membre ON Reservation.id_membre = Membre.id_membre
WHERE Membre.nom = "Alonso" AND Membre.prenom = "Fernando";
\end{lstlisting}
\end{reponse}
\resultat{\begin{tabular}{|c|}\hline COUNT(*) \\\hline 3 \\\hline\end{tabular}}
\methode
\begin{enumerate}
\item Le nom est dans \code{Membre}, les réservations dans \code{Reservation} : il faut les relier par une jointure sur \code{id\_membre}.
\item On garde les lignes où le nom est Alonso \textbf{et} le prénom Fernando.
\item \code{COUNT(*)} compte les lignes gardées : les réservations 2, 5 et 6, soit 3.
\end{enumerate}

\begin{erreur}
Oublier la condition \sql{ON} et écrire \code{FROM Reservation, Membre WHERE nom = "Alonso"} : chaque réservation est associée à chaque membre, et on compte 6 au lieu de 3 (sur la table du sujet). La condition d'égalité des clés est indispensable.
\end{erreur}

\textbf{b)} \rappel{les noms et prénoms des membres qui possèdent un appareil à raclette. \hfill\textit{(2 points)}}
\begin{reponse}
\begin{lstlisting}
SELECT Membre.nom, Membre.prenom
FROM Membre
JOIN Possede ON Membre.id_membre = Possede.id_membre
JOIN Objet ON Possede.id_objet = Objet.id_objet
WHERE Objet.description = "Appareil à raclette";
\end{lstlisting}
\end{reponse}
\resultat{\begin{tabular}{|l|l|}\hline nom & prenom \\\hline Ali & Mohamed \\ Alonso & Fernando \\\hline\end{tabular}}
\explication Le nom des membres est dans \code{Membre}, la description dans \code{Objet}, et le lien entre les deux est dans \code{Possede}. Il faut donc \textbf{deux} jointures : \code{Membre} avec \code{Possede} (par \code{id\_membre}), puis \code{Possede} avec \code{Objet} (par \code{id\_objet}). C'est exactement le chemin suivi à la main à la question 1~b, et on retrouve le même résultat.

\end{document}
